\chapter{Introduction}
\label{ch:introduction}

Corporate insolvencies are rising worldwide for the fifth consecutive year.
\textcite{allianztrade2025} expect a further increase of $5\%$ in 2026, taking
global business failures to a record level about a quarter above where they stood
before the pandemic and putting an estimated $2.1$ million jobs at risk. Behind
each of these failures stand creditors, employees, suppliers and banks who carry
the cost, which is why predicting bankruptcy early enough to react has been a
research question for decades. Machine learning has pushed this field to
remarkable reported numbers, with recent studies reaching accuracies of $98$--$99\%$
\parencite{pham2025, suresh2025, sinap2025}. Two things are troubling about these
results. First, they rest on synthetic oversampling, which balances the data by
fabricating bankruptcy events that never occurred. Second, they are produced by
models whose decisions remain largely unexplained, so even a correct prediction
gives an analyst little to work with.

This thesis takes the opposite route on both points. Using a widely studied
dataset of $6{,}819$ Taiwanese firms, of which only $3.2\%$ went bankrupt, it
trains exclusively on real observations. The imbalance is handled by
ClusterCentroids undersampling, which condenses the healthy majority into
representative real firms instead of inventing bankrupt ones. Three models of
increasing complexity, a logistic regression, a random forest and gradient
boosting, are compared under repeated stratified cross-validation and selected on
the $F_2$-score, because missing a bankruptcy costs far more than a false alarm.
The setup extends \textcite{wangliu2021}, who established this sampling and
evaluation framework on the same data but stopped at predictive performance.

The added layer is the actual contribution and leads to the research question of
this thesis, which financial features the fitted models rely on and how far the
answer depends on how one asks. Three importance methods, native importance,
permutation importance and SHAP, are applied to all three models. The short answer
is that no single ranking exists, as the methods diverge in ways their known
weaknesses predict, while a consistent signal remains, with every model placing a
measure of debt dependence at or near the top. Chapters~2 and~3 review the
literature and the data, Chapters~4 and~5 develop the methodology and the
importance methods and Chapters~6 to~8 report, discuss and conclude.